\documentclass{ximera}

\graphicspath{{./}{thePythagoreanTheorem/}{deMoivreSavesTheDay/}{complexNumbersFromDifferentAngles/}}

\usepackage{tikz}
\usepackage{tkz-euclide}
\usetkzobj{all}
\tikzstyle geometryDiagrams=[ultra thick,color=blue!50!black]
\newcommand{\tri}{\triangle}
\renewcommand{\l}{\ell}
\renewcommand{\P}{\mathcal{P}}
\newcommand{\R}{\mathbb{R}}
\newcommand{\Q}{\mathbb{Q}}

\newcommand{\Z}{\mathbb Z}

\renewcommand{\vec}{\mathbf}
\renewcommand{\d}{\,d}



%% Egyptian symbols

\usepackage{multido}
\newcommand{\egmil}[1]{\multido{\i=1+1}{#1}{\includegraphics[scale=.1]{egyptian/egypt_person.pdf}\hspace{0.5mm}}}
\newcommand{\eghuntho}[1]{\multido{\i=1+1}{#1}{\includegraphics[scale=.1]{egyptian/egypt_fish.pdf}\hspace{0.5mm}}}
\newcommand{\egtentho}[1]{\multido{\i=1+1}{#1}{\includegraphics[scale=.1]{egyptian/egypt_finger.pdf}\hspace{0.5mm}}}
\newcommand{\egtho}[1]{\multido{\i=1+1}{#1}{\includegraphics[scale=.1]{egyptian/egypt_lotus.pdf}\hspace{0.5mm}}}
\newcommand{\eghun}[1]{\multido{\i=1+1}{#1}{\includegraphics[scale=.1]{egyptian/egypt_scroll.pdf}\hspace{0.5mm}}}
\newcommand{\egten}[1]{\multido{\i=1+1}{#1}{\includegraphics[scale=.1]{egyptian/egypt_heel.pdf}\hspace{0.5mm}}}
\newcommand{\egone}[1]{\multido{\i=1+1}{#1}{\includegraphics[scale=.1]{egyptian/egypt_stroke.pdf}\hspace{0.5mm}}}
\newcommand{\egyptify}[7]{
 \multido{\i=1+1}{#1}{\includegraphics[scale=.1]{egyptian/egypt_person.pdf}\hspace{0.5mm}}
 \multido{\i=1+1}{#2}{\includegraphics[scale=.1]{egyptian/egypt_fish.pdf}\hspace{0.5mm}}
 \multido{\i=1+1}{#3}{\includegraphics[scale=.1]{egyptian/egypt_finger.pdf}\hspace{0.5mm}}
 \multido{\i=1+1}{#4}{\includegraphics[scale=.1]{egyptian/egypt_lotus.pdf}\hspace{0.5mm}}
 \multido{\i=1+1}{#5}{\includegraphics[scale=.1]{egyptian/egypt_scroll.pdf}\hspace{0.5mm}}
 \multido{\i=1+1}{#6}{\includegraphics[scale=.1]{egyptian/egypt_heel.pdf}\hspace{0.5mm}}
 \multido{\i=1+1}{#7}{\includegraphics[scale=.1]{egyptian/egypt_stroke.pdf}\hspace{0.5mm}}
 \hspace{.5mm}
}




%\prerequisites{algebra}
%\outcomes{arithmeticMean,geometricMean,harmonicMean}

\title{Pythagorean means}

\begin{document}
\begin{abstract}
In this activity we explore the three different means of the
  ancient Greeks.
\end{abstract}
\maketitle


\subsection*{The arithmetic mean}

The \textit{arithmetic mean} is the good-old mean that we are all
familiar with.

\begin{question}
What is the mean that we are all familiar with? Explain how to compute
the mean of $a_1,a_2,\dots,a_n$. Give some examples.
\end{question}

\subsection*{The geometric mean}

The \textit{geometric mean} is a bit different. The geometric mean of $a_1,a_2,\dots, a_n$ is given by:
\[
\left(\prod_{i=1}^n a_i\right)^{1/n}
\]

\begin{question}
Explain an analogy between the arithmetic mean and the geometric mean.
\end{question}

\begin{question}
Can you explain the geometric mean in terms of geometry? First do it
for 2 numbers. Next do it for three.
\end{question}

\subsection*{The harmonic mean}

The \textit{harmonic mean} might be the most mysterious of all. The
harmonic mean of $a_1,a_2,\dots, a_n$ is given by:
\[
\frac{n}{\sum_{i=1}^n \frac{1}{a_i}}
\]

\begin{exploration}
Can you find a connection between the harmonic mean and music?
\end{exploration}


\begin{question}
In the United States, the fuel efficiency of a car is usually given in
the units:
\[
\frac{\text{miles}}{\text{gallon}}
\]
However, in Europe, the fuel efficiency of a car is usually given in
the units:
\[
\frac{\text{liters}}{100 \mathrm{km}}
\]
Give some examples of fuel efficiency (both efficient and
inefficient) with each set of units.
\end{question}


\begin{question}
Now suppose that a car gets $60\frac{\text{miles}}{\text{gallon}}$ and
another car gets $20\frac{\text{miles}}{\text{gallon}}$. What is the
average fuel efficiency?
\end{question}


\begin{question}
Now suppose that a car gets $4\frac{\text{liters}}{100 \mathrm{km}} $
and another car gets $20\frac{\text{liters}}{100 \mathrm{km}}$. What
is the average fuel efficiency?
\end{question}


\begin{exploration}
Compare your answers to the last two questions. Something fishy is
going on, what is it?
\end{exploration}
\end{document}
